\section{Detailed one-step ledger estimates}\label{app:one-step}

This appendix records the estimates used in \cref{sec:residual,sec:ledger} in
a form independent of the exploratory calculation log.  Throughout, $C$ is
fixed and sufficiently large, $w>\omega_0$, and the target set exposed after
column $i$ is $V_i=\{i+1,\ldots,r\}$, of size $k=r-i$.

\subsection{The Hessian on the entire cutoff box}

Recall
\[
 F_i(b)=\sum_{j\in V_i}\log\Phi(-b_j)
 -\frac1d\sum_{\{j,q\}\subset V_i}T_{jq}m(b_j)m(b_q).
\]
Since $(\log\Phi(-x))'=-m(x)$, the entries of the Hessian are
\begin{align}
 \partial_{jj}^2F_i(b)
 &=-m'(b_j)-\frac{m''(b_j)}d
   \sum_{q\in V_i\setminus\{j\}}T_{jq}m(b_q), \tag{A.1}\\
 \partial_{jq}^2F_i(b)
 &=-\frac{T_{jq}}d m'(b_j)m'(b_q),\qquad j\ne q. \tag{A.2}
\end{align}
The inverse Mills ratio satisfies $m''\ge0$.  Hence the interaction term in
(A.1) can only strengthen the diagonal of $-\nabla^2F_i$.  On
$[c-w,c+w]^{V_i}$, monotonicity gives
\[
 m'(b_j)\ge j_-(w),\qquad |m'(b_j)m'(b_q)|\le j_+(w)^2.
\tag{A.3}
\]

The residual recursion and $r/\sqrt d\le1/D$ imply
\begin{align}
 T_{\max}&\le \frac{t_0+\theta r}{1-\rho}, \tag{A.4}\\
 \frac1d\max_j\sum_{q\ne j}T_{jq}
 &\le \frac{m_0+\theta/D}{D(1-\rho)}. \tag{A.5}
\end{align}
Combining (A.1)--(A.5) with Gershgorin's theorem proves
\[
 -\nabla^2F_i(b)\succeq
 \left[j_-(w)-j_+(w)^2
 \frac{m_0+\theta/D}{D(1-\rho)}\right]I
 =\mu_w I. \tag{A.6}
\]
This is a full-box statement, not an expansion valid only at the central
cutoff.

\subsection{Boundary merge with a controlled residual}

The gradient at $c\mathbf1$ is
\[
 \partial_jF_i(c\mathbf1)
 =-m_0-\frac{m_0m'_0}{d}
   \sum_{q\in V_i\setminus\{j\}}T_{jq}.
\tag{A.7}
\]
The defining recursion
\[
 T_{ij}=m_0\sqrt d+\frac{m_0m'_0}{\sqrt d}
 \sum_{q\in V_i\setminus\{j\}}T_{jq}+e_{ij}
\tag{A.8}
\]
therefore gives
\[
 \partial_jF_i(c\mathbf1)=-\frac{T_{ij}-e_{ij}}{\sqrt d}.
\tag{A.9}
\]
For $\delta_j=b_j-c=\sqrt d\,h_{ij}+\varepsilon_{ij}$, with
$|\varepsilon_{ij}|\le d^{-1/5}$, the tangent inequality from (A.6) gives
\begin{align}
 F_i(b)
 &\le F_i(c\mathbf1)
 -\sum_{j\in V_i}T_{ij}h_{ij}
 +\sum_{j\in V_i}\left(
   \frac{e_{ij}\delta_j}{\sqrt d}-\frac{\mu_w\delta_j^2}{2}\right)
 +\frac{d^{-1/5}}{\sqrt d}\sum_{j\in V_i}T_{ij} \notag\\
 &\le F_i(c\mathbf1)-\sum_{j\in V_i}T_{ij}h_{ij}
 +\sum_{j\in V_i}\frac{e_{ij}^2}{2\mu_wd}
 +\frac{d^{-1/5}}{\sqrt d}\sum_{j\in V_i}T_{ij}. \tag{A.10}
\end{align}
The second line is the scalar square completion
$xy-\mu_wy^2/2\le x^2/(2\mu_w)$.  Thus every boundary edge merges with
its full coefficient $T_{ij}$; only the explicit residual is paid.

\subsection{Mass identities}

Let $\Delta_{ij}=T_{ij}-t_0$.  Backward induction in the first endpoint
gives, for every $j<q$,
\[
 A_0(r-j-1)+\theta(j-1)
 \le \Delta_{jq}
 \le \frac{A_0+\rho\theta}{1-\rho}(r-j-1)+\theta(j-1).
\tag{A.11}
\]
For $S_i$ and $R_i$ as in \cref{sec:ledger}, this implies
\[
 S_i\ge a_-(r-2)\binom{k}{2},\qquad
 R_i\le a_+(r-2)(k-1),\qquad
 \frac{R_i}{S_i}\le\frac{2a_+}{a_-k}. \tag{A.12}
\]

The exact accumulated sums are obtained by exchanging the order of
summation.  The future-degree part of (A.11) counts each ordered quadruple
twice, while the residual part contributes
\[
 K_r:=\sum_{i=1}^{r-1}(r-i)(i-1)^2
 =\binom r3+2\binom r4. \tag{A.13}
\]
Consequently
\[
 \sum_{i=1}^{r-1}S_i\ge2A_0\binom r4+\theta K_r,
 \qquad
 \sum_{1\le i<j\le r}e_{ij}^2=\theta^2K_r. \tag{A.14}
\]
No asymptotic replacement is used in (A.13)--(A.14).

\subsection{Three-factor H\"older comparison}

Fix a red-admissible \(\mathcal F_{i-1}\)-history for target \(r\), condition
further on \(A_i^R\), and write
\[
 X_j=-\frac{m(b_j)}{\sqrt d}+\frac{\zeta_j}{\sqrt d},\qquad
 u_j^0=\frac{t_0}{d}\sum_{q\ne j}m(b_q),\qquad
 \delta u_j=\frac1d\sum_{q\ne j}\Delta_{jq}m(b_q).
\tag{A.15}
\]
The variables \(\zeta_j\) are conditionally independent, centered and
\(1\)-subgaussian.  The centered exponent is exactly
\[
 \sum_j(u_j^0+\delta u_j)\zeta_j
 -\frac{t_0}{d}\sum_{j<q}\zeta_j\zeta_q
 -\frac1d\sum_{j<q}\Delta_{jq}\zeta_j\zeta_q.
\tag{A.16}
\]
Use conjugate exponents
\[
 Q_0=\frac d{4t_0k},\qquad Q_1=\frac d{4R_i},\qquad
 P^{-1}=1-Q_0^{-1}-Q_1^{-1}. \tag{A.17}
\]
They are all greater than one, with \(P\le2\), because
\[
 \frac{4m_0}{D}+\frac{4a_+}{D^2}<\frac12. \tag{A.18}
\]
H\"older gives one linear-CGF factor and two quadratic factors.  The
uniform \(t_0\) factor is the frozen HMS factor.  For the increment, use
\[
 -\sum_{j<q}\Delta_{jq}\zeta_j\zeta_q
 \le\frac12\sum_jR_j\zeta_j^2,\qquad
 \frac{Q_1R_j}{2d}\le\frac18.
\tag{A.19}
\]
The square-MGF estimate
\(\log\mathbb E e^{x\zeta_j^2}\le(16/3)x\) for \(0\le x\le1/8\)
therefore gives
\[
 \frac1{Q_1}\sum_j\frac{16}{3}\frac{Q_1R_j}{2d}
 =\frac{16}{3}\frac{S_i}{d}.
\tag{A.20}
\]

For the linear factor, let \(P^0=P^0_{k,d}\).  The source exact-CGF deficit
on this identical history is
\[
 \mathfrak d_i(b)=\sum_j\left[(u_j^0)^2
 -\frac1{P^0}K_{b_j}(P^0u_j^0)\right]
 \ge\mathsf d_{k,w,d}.
\tag{A.21}
\]
Thus it remains only to bound the change from
\((P^0,u^0)\) to \((P,u^0+\delta u)\).  Since \(0\le K_b''\le1\),
\begin{align}
 \frac{K_b(P(u+\delta))-K_b(Pu)}P
 &\le2u\delta+\delta^2,\tag{A.22}\\
 0\le\frac{\partial}{\partial P}\frac{K_b(Pu)}P
 &\le u^2.\tag{A.23}
\end{align}
Furthermore
\[
 P-P^0\le\frac{16R_i}{d},\quad
 \sum_j\delta u_j\le\frac{2M_wS_i}{d},\quad
 \max_j\delta u_j\le\frac{M_wR_i}{d},
\tag{A.24}
\]
and
\[
 \sum_j(u_j^0)^2\le
 \frac{2m_0^2M_w^2}{a_-}\frac{S_i}{d}.
\tag{A.25}
\]
Substitution in (A.22)--(A.24) gives, in order,
\[
 \frac{4m_0M_w^2}{D}\frac{S_i}{d},\qquad
 \frac{2M_w^2a_+}{D^2}\frac{S_i}{d},\qquad
 \frac{32m_0^2M_w^2a_+}{a_-D^2}\frac{S_i}{d}.
\tag{A.26}
\]
These are precisely the three summands of
\(\widehat L_w(C)S_i/d\).  Combining them with (A.20) and the central
deterministic gain \(-m_0^2S_i/d\) gives
\[
 -\left(m_0^2-\frac{16}{3}-\widehat L_w(C)\right)\frac{S_i}{d}
 =-\widehat B_w(C)\frac{S_i}{d}. \tag{A.27}
\]
Equation (A.21) is charged once, (A.27) pays only the new
\(\Delta=T-t_0\) factor and the change of exponent, and (A.10) pays the
residual gradient.  This is the promised same-history separation.

\subsection{One-step conclusion and extraction}

Combining (A.10), (A.21), and (A.27) proves (2.37) with the value
\(\lambda_i\) in (6.2).  The weighted reverse-propagation
\cref{lem:weighted-reverse}, followed by (A.14), then gives the empty-state bound
\[
 P^*_{R,r}\le\mathcal H_C(r)
 \exp\!\left[-\mathcal D_{C,w,d}(r)
 -\frac1d\left\{2\widehat B_wA_0\binom r4+E_wK_r\right\}
 +J_{r,d}\right]. \tag{A.28}
\]
where
\[
 0\le J_{r,d}\le
 \frac{m_0+\theta/D}{1-\rho}d^{-1/5}\binom r2=o_C(r^2).
 \tag{A.29}
\]
If extraction deletes $u$ vertices from a retained set of size $\ell$,
then the old exact-CGF restoration cost is at most
\[
 \mathcal D_{C,w,d}(\ell)-\mathcal D_{C,w,d}(\ell-u)
 \le\frac{m_0^4}{D^2}\ell u\le\ell u, \tag{A.30}
\]
and the loss in the new deterministic exponent is at most
\[
 \frac1d\left[
 2\widehat B_wA_0u\binom\ell3
 +E_w\left(u\binom\ell2+2u\binom\ell3\right)
 \right]
 \le\left[\frac{\widehat B_wA_0+E_w}{3D^2}
 +\frac{E_w}{2D^2\ell}\right]\ell u.
 \tag{A.31}
\]
For sufficiently large \(C\) and then sufficiently large \(\ell\), the
bracket in (A.31) is
\[
 O(1/\log C)+O_C(\ell^{-1})<1.
\]
Together
with the source-envelope restoration
\([\log(1/p)+1]\ell u\), the total cost is at most
\([\log(1/p)+3]\ell u\).  The pinned failure factor therefore leaves
\(\eta_C=10\log(10/p)-\log(1/p)-3>0\), and
\[
 \sum_{u=1}^{\ell}\binom\ell u e^{-\eta_C\ell u}
 =(1+e^{-\eta_C\ell})^\ell-1=o(1). \tag{A.32}
\]
This proves that (A.28) survives perfect-sequence extraction with its entire
old and new exponent intact.
